\section{Cartesian Feedback Control and Motion Execution}
\label{sec:feedback_control}

\subsection{Control Problem and Baseline}

The numerical inverse-kinematic solver produces joint angles that satisfy the
model-based convergence criterion. After conversion to integer ticks, however,
a single command does not guarantee that the servos reach these targets
exactly. Each STS3215 controls only its own motor position
\cite{feetech_protocol}, not the Cartesian result of the coupled motion.

The one-shot baseline therefore read the four joint positions, solved the
complete IK problem, and sent all targets in one synchronous write. Motion was
considered settled when every encoder changed by at most one tick for three
cycles; a \(20\,\mathrm{s}\) timeout bounded failures. The measured ticks were
then evaluated through forward kinematics as

\begin{equation}
    \hat{\mathbf{p}}(k)
    =
    \mathbf{p}\!\left(\mathbf{q}_{\mathrm{meas}}(k)\right).
    \label{eq:estimated_tcp}
\end{equation}

This estimate is not an independently measured physical tool center point
(TCP). Encoder feedback and the kinematic model cannot reveal backlash, link
deflection, geometric calibration errors, or displacement of the physical TCP.

\subsection{Resolved-Rate P Controller}

Resolved-rate control converts Cartesian error into coordinated joint motion
through the manipulator Jacobian \cite{whitney1969resolved}. Each cycle begins
with a synchronous read of the four active joints. The Cartesian error is

\begin{equation}
    \mathbf{e}_{p}(k)
    =
    \mathbf{p}_{\mathrm{target}}-\hat{\mathbf{p}}(k).
    \label{eq:resolved_rate_error}
\end{equation}

The outer controller applies proportional feedback with a speed limit:

\begin{equation}
    \mathbf{v}_{c}(k)
    =
    \operatorname{sat}_{v_{\max}}
    \left(K_{P,c}\mathbf{e}_{p}(k)\right),
    \qquad
    K_{P,c}=2\,\mathrm{s}^{-1},
    \quad
    v_{\max}=0.1\,\mathrm{m/s}.
    \label{eq:cartesian_p_velocity}
\end{equation}

Using the measured cycle time, rather than assuming the nominal
\(20\,\mathrm{ms}\) period, gives

\begin{equation}
    \Delta\mathbf{p}(k)=\mathbf{v}_{c}(k)\Delta t_k,
    \qquad
    \mathbf{p}_{\mathrm{next}}(k)
    =
    \hat{\mathbf{p}}(k)+\Delta\mathbf{p}(k).
    \label{eq:resolved_rate_cartesian_step}
\end{equation}

Blocking communication and telemetry can vary the actual interval. Including
\(\Delta t_k\) keeps displacement proportional to elapsed time rather than the
number of controller calls. At the nominal period, for example, an unsaturated
\(20\,\mathrm{mm}\) error produces a \(0.8\,\mathrm{mm}\) Cartesian step.

One damped least-squares update maps this displacement into joint space:

\begin{equation}
    \Delta\mathbf{q}(k)
    \approx
    \mathbf{J}^{+}(\mathbf{q})\mathbf{v}_{c}(k)\Delta t_k.
    \label{eq:resolved_rate_joint_step}
\end{equation}

The floating-point joint command accumulates these increments so that sub-tick
corrections are not lost. After joint-limit checks, all targets are converted
to ticks and sent synchronously. The loop terminates when
\(\lVert\mathbf{e}_{p}\rVert_2\leq1\,\mathrm{mm}\), an abort is requested, or
the timeout expires. Although the actuators receive joint-position targets,
the feedback error and controller output are Cartesian; this distinguishes the
method from independent joint-space control.

The accumulated command can temporarily lead the measured joint positions if
the servos cannot follow each increment immediately. This effect may contribute
to endpoint oscillation, but it is not integral windup because the implemented
outer controller contains no integral state.

\subsection{Selection of the Control Terms}

For errors above \(v_{\max}/K_{P,c}=50\,\mathrm{mm}\), speed saturation
dominates. Increasing \(K_{P,c}\) to \(4\,\mathrm{s}^{-1}\) delayed
deceleration and produced more endpoint correction. A filtered
derivative-on-measurement term was tested to avoid derivative kick and reduce
the influence of encoder quantization. It did not improve settling and produced
more corrective motion near the waypoint. Because no independent TCP
measurement was available, this is reported as observed oscillation rather
than quantified physical overshoot. An integral term offered no demonstrated
benefit and would require anti-windup during Cartesian or joint saturation.
The final implementation therefore retained the simpler P controller
\cite{astrom2008feedback}.

\subsection{Experimental Result and Limitation}

The baseline and controlled runs used the same start pose, four square
waypoints, speed, acceleration, timeout, model, and joint limits. No E-Skin
pause or communication fault occurred. Table~\ref{tab:controller_endpoint_error}
shows terminal TCP errors reconstructed from servo feedback.

\begin{table}[!h]
    \centering
    \caption{Model-based terminal TCP error for one hardware run per method.}
    \label{tab:controller_endpoint_error}
    \begin{tabular}{lrr}
        \hline
        Waypoint & No outer controller & Resolved rate \\
        \hline
        Upper left  & \(16.44\,\mathrm{mm}\) & \(0.93\,\mathrm{mm}\) \\
        Upper right & \(12.60\,\mathrm{mm}\) & \(0.84\,\mathrm{mm}\) \\
        Lower right & \(4.21\,\mathrm{mm}\)  & \(0.90\,\mathrm{mm}\) \\
        Lower left  & \(7.28\,\mathrm{mm}\)  & \(0.81\,\mathrm{mm}\) \\
        \hline
        Mean        & \(10.13\,\mathrm{mm}\) & \(0.87\,\mathrm{mm}\) \\
        Maximum     & \(16.44\,\mathrm{mm}\) & \(0.93\,\mathrm{mm}\) \\
        \hline
    \end{tabular}
\end{table}

The controller remained below \(1\,\mathrm{mm}\) at every waypoint and reduced
the mean model-based error
by approximately \(91.4\,\%\), or a factor of \(11.6\). This result justified
its selection for the final demonstration. However, one run per method does
not establish physical accuracy or repeatability. Both evaluation values
depend on the same forward model; independent metrology and repeated trials
would be required for those claims.
